\documentclass[11pt]{article}

\usepackage[margin=1in]{geometry}
\usepackage{amsmath, amssymb, amsthm}
\usepackage{enumitem}

\newcommand{\Z}{\mathbb{Z}}
\newcommand{\R}{\mathbb{R}}

\pagestyle{empty}

\begin{document}

\begin{center}
  {\Large\bfseries Weekly Assignment: Combinations and the Binomial Theorem (Epp Ch.\ 9)}
\end{center}

\bigskip

\noindent
Complete each of the following problems.  Show all work and justify your answers.

\bigskip

%% ─────────────────────────────────────────────
\noindent\textbf{Problem 1} (Epp 9.5 \#2).
\begin{enumerate}[label=\alph*.]
  \item List all 3-combinations for the set $\{x_1, x_2, x_3, x_4, x_5\}$.
        Deduce the value of $\displaystyle\binom{5}{3}$.
  \item List all unordered selections of two elements from the set
        $\{x_1, x_2, x_3, x_4, x_5, x_6\}$.
        Deduce the value of $\displaystyle\binom{6}{2}$.
\end{enumerate}

\bigskip

%% ─────────────────────────────────────────────
\noindent\textbf{Problem 2} (Epp 9.5 \#4). \;
Write an equation relating $P(8, 3)$ and $\displaystyle\binom{8}{3}$.

\bigskip

%% ─────────────────────────────────────────────
\noindent\textbf{Problem 3} (Epp 9.5 \#7). \;
A computer programming team has 13 members.
\begin{enumerate}[label=\alph*.]
  \item How many ways can a group of seven be chosen to work on a project?
  \item Suppose seven team members are women and six are men.
    \begin{enumerate}[label=(\roman*)]
      \item How many groups of seven can be chosen that contain four women
            and three men?
      \item How many groups of seven can be chosen that contain at least
            one man?
      \item How many groups of seven can be chosen that contain at most
            three women?
    \end{enumerate}
  \item Suppose two team members refuse to work together on projects.
        How many groups of seven can be chosen to work on a project?
  \item Suppose two team members insist on either working together or
        not at all on projects.
        How many groups of seven can be chosen to work on a project?
\end{enumerate}

\bigskip

%% ─────────────────────────────────────────────
\noindent\textbf{Problem 4} (Epp 9.6 \#4). \;
A camera shop stocks eight different types of batteries, one of which is
type~A76.  Assume there are at least 30 batteries of each type.
\begin{enumerate}[label=\alph*.]
  \item How many ways can a total inventory of 30 batteries be distributed
        among the eight different types?
  \item How many ways can a total inventory of 30 batteries be distributed
        among the eight different types if the inventory must include at
        least four A76 batteries?
  \item How many ways can a total inventory of 30 batteries be distributed
        among the eight different types if the inventory includes at most
        three A76 batteries?
\end{enumerate}

\bigskip

%% ─────────────────────────────────────────────
\noindent\textbf{Problem 5} (Epp 9.7 \#12). \;
Use Pascal's formula repeatedly to derive a formula for
$\displaystyle\binom{n+3}{r}$ in terms of values of $\displaystyle\binom{n}{k}$
with $k \le r$.
(Assume $n$ and $r$ are integers with $n \ge r \ge 3$.)

\end{document}
